% ============================================
% Chapter 3: Algorithm Design
% ============================================

\chapter{Algorithm Design}

This chapter details the design philosophy and implementation of each dispatching strategy. All algorithms share the same snapshot-to-command interface and a common utility toolbox, making them easy to compare, extend, and benchmark under identical simulation conditions.

\section{Unified Scheduling Framework}

\subsection{Snapshot and Command}

Every scheduling round begins with a read-only snapshot captured from the data layer, containing:
\begin{itemize}[leftmargin=4em]
    \item All vehicles: current positions, battery, payload, status, and on-board tasks.
    \item All tasks with their lifecycle states.
    \item All charging stations: positions, capacity, queue length, and load pressure.
    \item The current simulation time $t_{\text{sim}}$.
    \item Cached shortest-path queries over the road network.
\end{itemize}

The algorithm returns a list of commands. Each command specifies one vehicle's \emph{next step}:
\begin{itemize}[leftmargin=4em]
    \item \textbf{Deliver:} go to a task location (optionally loading a batch at the warehouse).
    \item \textbf{Charge:} go to a charging station.
    \item \textbf{Return:} go back to the warehouse.
    \item \textbf{Go to node:} move to an arbitrary coordinate (e.g., a relay meeting point).
    \item \textbf{Handoff:} transfer cargo to another co-located vehicle.
    \item \textbf{Idle:} stay put.
\end{itemize}

Multi-stop delivery is achieved by repeated snapshot-to-decision cycles: after each leg completes, the vehicle becomes idle and the scheduler is invoked again.

\subsection{Shared Utility Toolbox}

A shared utility module centralizes logic that every strategy would otherwise re-implement. Table~\ref{tab:utils} summarizes the main capability groups; developers can use these defaults as-is or override only the task-selection step:

\begin{table}[H]
    \centering
    \caption{Shared Utility Capability Groups}
    \label{tab:utils}
    \reporttableset{0.98\textwidth}
    \begin{reporttabular}{>{\raggedright\arraybackslash}l >{\raggedright\arraybackslash}l}
        \toprule
        \textbf{Category} & \textbf{Capabilities} \\
        \midrule
        Query \& filter & Road-network distance, payload-feasible task filtering, nearest task/station lookup \\
        Battery \& charge & Reachability checks, chain-distance estimation with mid-route recharge, low-battery detection, optimal station selection \\
        Route ordering & Greedy nearest-neighbor chaining, MST-based visit ordering \\
        Command construction & Standard builders for deliver, charge, return, handoff, and idle actions \\
        Decision templates & Warehouse and en-route decision flows with built-in battery safety \\
        Relay helpers & Path midpoint calculation, cargo-transfer feasibility checks \\
        \bottomrule
    \end{reporttabular}
\end{table}

To implement a custom scheduler, one typically overrides only the task-picking logic; battery checks, charging fallback, route ordering, and command construction are handled automatically by the shared templates.

\subsection{Battery Pre-Check and Mid-Route Recharge}

Before dispatching a vehicle, the system verifies battery feasibility for the entire planned chain, simulating mid-route charging detours when direct travel is insufficient. Algorithm~\ref{algo:battery_check} outlines this chain-distance estimation procedure.

\begin{algorithm}[H]
    \caption{Battery Pre-check with Mid-Route Recharge}
    \label{algo:battery_check}
    \begin{algorithmic}[1]
        \Function{EstimateChainDistance}{vehicle, waypoints, snapshot}
            \State $cur \gets$ vehicle position,\ $battery \gets$ vehicle battery level
            \State $totalDist \gets 0$,\ $\eta \gets$ unit energy consumption,\ $\gamma \gets$ safety margin (1.15)
            \For{each $target$ in waypoints}
                \Loop
                    \State $need \gets$ energy to reach $target$ (+ station buffer if required) $\times \gamma$
                    \If{$battery \ge need$}
                        \State Advance $cur$ to $target$; update $totalDist$ and $battery$; \textbf{break}
                    \EndIf
                    \Comment{Insufficient: detour to best reachable transit station, recharge to 90\%}
                    \State $station \gets$ best transit station from $cur$ toward $target$
                    \If{no reachable station} \Return infeasible \EndIf
                    \State Detour to $station$; reset $battery$ to charge target; continue loop
                \EndLoop
            \EndFor
            \State \Return $totalDist$
        \EndFunction
    \end{algorithmic}
\end{algorithm}

\subsection{Decision Templates}

Most per-vehicle strategies follow a two-context pattern: warehouse decisions and en-route decisions, as shown in Algorithm~\ref{algo:template}.

\begin{algorithm}[H]
    \caption{Unified Per-Vehicle Decision Template}
    \label{algo:template}
    \begin{algorithmic}[1]
        \Require Snapshot $snap$, custom task-picking rule
        \Ensure List of commands
        \State $commands \gets$ empty list
        \Comment{En-route vehicles: threshold charge, next undelivered task, or return}
        \For{each en-route idle vehicle $v$}
            \State Append en-route decision for $v$
        \EndFor
        \Comment{Warehouse vehicles: pick batch via strategy-specific logic}
        \State $claimed \gets$ empty set
        \For{each warehouse idle vehicle $v$}
            \State $cmd \gets$ warehouse decision for $v$ using the custom task-picking rule
            \If{$cmd$ assigns new tasks} \State $claimed$.update(assigned task ids) \EndIf
            \State $commands$.append($cmd$)
        \EndFor
        \State \Return $commands$
    \end{algorithmic}
\end{algorithm}

At the warehouse, the template sorts the picked batch via greedy route ordering, checks full-chain battery feasibility, and issues either a deliver command (loading all tasks) or a transit charge command if the first hop requires recharging first.

\section{Greedy Baseline Strategies}

Four greedy baselines differ only in how they sort candidate tasks before greedily building a maximal feasible prefix batch. All share the same warehouse and en-route decision templates.

\subsection{Nearest Task First}

Sorts unclaimed tasks by distance from the vehicle; greedily adds tasks in nearest-first order while the chain remains payload- and battery-feasible. The simplest and fastest baseline, suitable when response time matters.

\subsection{Heaviest Task First}

Sorts by descending cargo weight (tie-break: distance). Assigns heavy cargo early to prevent payload-constrained vehicles from being stuck with only light tasks later.

\subsection{Earliest Deadline First}

Sorts by ascending deadline (tie-break: distance). A classic real-time scheduling heuristic (EDF) that prioritizes time-sensitive deliveries.

\subsection{Priority First}

Sorts by descending priority (tie-break: distance). Suitable when VIP or urgent orders must be served first regardless of distance.

All four follow the batch-building pattern shown in Algorithm~\ref{algo:greedy_batch}.

\begin{algorithm}[H]
    \caption{Greedy Maximal Feasible Batch (shared by all four baselines)}
    \label{algo:greedy_batch}
    \begin{algorithmic}[1]
        \Function{PickBatch}{vehicle, tasks, snapshot, sort rule}
            \State $candidates \gets$ payload-feasible tasks sorted by the given rule
            \State $chosen \gets [\ ]$
            \For{each task $t$ in $candidates$}
                \State $ordered \gets$ greedy route order for $chosen \cup \{t\}$
                \If{chain distance infeasible} \textbf{break} \EndIf
                \State $chosen \gets chosen \cup \{t\}$
            \EndFor
            \State \Return $chosen$
        \EndFunction
    \end{algorithmic}
\end{algorithm}

\section{Search-Based Batch Loaders}

These strategies replace the greedy prefix rule with explicit search over feasible task subsets, using a unified batch chain score as the objective.

\subsection{DFS Score Search}

Enumerates all \textbf{maximal feasible batches}---subsets where no additional single task can be added without violating payload or battery constraints---via depth-first search with payload pruning. Selects the batch with the highest predicted chain score. Guarantees the optimal batch per vehicle per decision round among maximal feasible sets, at exponential worst-case cost (mitigated by small pending queues in practice).

\subsection{SA Score Search}

Uses simulated annealing to approximate the DFS optimum over the same maximal-feasible batch space. Neighbor moves add, remove, or swap tasks within the candidate pool; the cooling schedule balances exploration and exploitation. Faster than DFS when pending task count is large, with a reheat mechanism to escape local optima.

\subsection{RL Batch}

A hybrid approach combining heuristic pruning with reinforcement learning:
\begin{itemize}[leftmargin=4em]
    \item \textbf{Candidate pool:} Build a nearest-first greedy prefix, then cap the pool size (default 7 tasks).
    \item \textbf{Batch selection:} Run DFS score search within the capped pool to pick the best maximal batch.
    \item \textbf{Online learning:} Maintain a Q-table over state--action pairs where state encodes nearest-count and battery tier; update Q-values from observed batch rewards after delivery completion ($\epsilon$-greedy exploration with decay).
\end{itemize}

This design keeps per-decision latency bounded while allowing the policy to adapt across repeated simulations.

\section{Multi-Vehicle Coordination Strategies}

These strategies go beyond independent per-vehicle decisions and explicitly coordinate fleet-level task allocation.

\subsection{Regional Planning}

When pending tasks $\ge 5$:
\begin{enumerate}[leftmargin=4em]
    \item Cluster tasks via K-means ($k=5$) into geographic \emph{packages}, sorted by total weight.
    \item Idle warehouse vehicles \emph{claim} packages: a vehicle that can carry the entire package competes on predicted batch score; if no vehicle fits, defer if a returning empty vehicle could take it, otherwise assign the best-scoring feasible subset.
    \item Remaining idle vehicles pick up leftover packages, prioritizing the vehicle farthest from the warehouse.
\end{enumerate}
For fewer than 5 pending tasks, the strategy degrades to Nearest Task First.

\subsection{Cluster Auction MAS}

Models each idle warehouse vehicle as an independent \textbf{agent}. Tasks are clustered into lots (K-means packages or single-task lots when pending $< 5$). Agents submit bids:
\begin{equation}
    \text{bid}(v, lot) = \text{predicted score}(v, lot) - \alpha \times d(v, \text{lot centroid})
\end{equation}
A central coordinator runs a \textbf{sequential auction}: lots are offered in descending weight order; the highest feasible full-package bid wins (one lot per vehicle per round). If no full-package bid exists, partial-batch bids are accepted. This multi-agent design distributes decision-making while avoiding the combinatorial explosion of a single global optimizer.

\subsection{Relay Handoff}

Enables \textbf{in-route cargo exchange} between nearby vehicles:
\begin{enumerate}[leftmargin=4em]
    \item Warehouse batching follows the Nearest Task First rule.
    \item Each round, scan all non-stranded vehicle pairs outside the warehouse whose road-network distance is below a proximity threshold.
    \item Merge on-board tasks, K-means split into two clusters; the lighter-load vehicle takes the lighter cluster (ascending weight order) and the other takes the remainder, subject to payload feasibility.
    \item Vehicles meet at the path midpoint; the first arrival waits; upon co-location, a handoff command transfers cargo.
    \item Each vehicle pair may exchange at most once per outbound-return cycle to prevent oscillation.
\end{enumerate}

This strategy explores cooperative delivery patterns that single-vehicle algorithms cannot achieve, such as rebalancing cargo mid-route when two vehicles converge.

\section{Algorithm Complexity Summary}

Table~\ref{tab:complexity} summarizes the time complexity of all ten algorithms, assuming $V$ idle vehicles and $T$ pending tasks per decision round.

\begin{table}[H]
    \centering
    \caption{Algorithm Complexity Comparison}
    \label{tab:complexity}
    \reporttableset{0.98\textwidth}
    \begin{reporttabular}{>{\raggedright\arraybackslash}l c >{\raggedright\arraybackslash}l}
        \toprule
        \textbf{Algorithm} & \textbf{Time Complexity} & \textbf{Characteristics} \\
        \midrule
        Nearest Task First       & $O(V \cdot T \log T)$ & Simplest greedy baseline \\
        Heaviest Task First      & $O(V \cdot T \log T)$ & Weight-prioritized greedy \\
        Earliest Deadline First  & $O(V \cdot T \log T)$ & Deadline-prioritized (EDF) \\
        Priority First           & $O(V \cdot T \log T)$ & Priority-prioritized greedy \\
        DFS Score Search         & $O(2^{m})$ & Optimal maximal batch; $m$ = pool size \\
        SA Score Search          & $O(N_{\text{iter}} \cdot m)$ & SA approximation of DFS \\
        RL Batch                 & $O(2^{k} + |Q|)$ & Capped pool ($k \le 7$) + Q-table lookup \\
        Regional Planning        & $O(k \cdot V \cdot T)$ & K-means + per-package scoring \\
        Cluster Auction MAS      & $O(k \cdot V \cdot T)$ & K-means + sequential auction \\
        Relay Handoff            & $O(V^2 \cdot T)$ & Pair scan + handoff negotiation \\
        \bottomrule
    \end{reporttabular}
\end{table}

\section{Extending with Custom Algorithms}

Adding a new strategy requires three steps:
\begin{enumerate}[leftmargin=4em]
    \item Implement a new scheduler class that reads the snapshot and returns commands.
    \item Register the class in the dynamic scheduler registry.
    \item Select the new strategy name in the YAML configuration file.
\end{enumerate}

For most custom heuristics, overriding only the task-picking logic inside the warehouse decision template is sufficient; the shared utilities handle battery safety, charging fallback, route ordering, and command construction automatically.
